\section{Diseño experimental: convertir "efectivo" en evidencia reproducible}

\subsection{Diseño de la matriz experimental}
El curso propone dividir los experimentos con Agentes en dos grupos: "variables de capacidad" y "variables de sistema": las variables de capacidad incluyen el reasoning budget, la estrategia de memory y la granularidad de planning; las variables de sistema incluyen la latencia de las herramientas, la tasa de aciertos de caché, la estrategia de concurrencia y las reglas de reintento ante fallos. Solo cuando ambos grupos de variables se controlan al mismo tiempo, las conclusiones del experimento tienen poder explicativo.

\begin{knowledgebox}{Unidad experimental mínima reproducible}
\begin{itemize}
    \item Fijar el conjunto de tareas y la versión de los datos;
    \item Fijar la interfaz de herramientas y los límites de permisos;
    \item Fijar la versión del modelo y el parámetro de temperatura;
    \item Registrar la trayectoria completa y los estados intermedios.
\end{itemize}
\end{knowledgebox}

\begin{table}[H]
\centering
\begin{tabular}{p{3.1cm}p{3.3cm}p{2.7cm}p{2.2cm}}
\toprule
\textbf{Elemento experimental} & \textbf{Configuración de control} & \textbf{Métrica principal} & \textbf{Métrica secundaria} \\
\midrule
Reasoning budget & 2k / 6k / 12k tokens & Tasa de éxito de la tarea & Costo por tarea \\
Memory policy & Sin memoria / memoria por recuperación / memoria jerárquica & Tasa de finalización de tareas largas & Duración de la recuperación ante errores \\
Planning policy & Plan estático / replanificación dinámica & Número promedio de pasos & Tasa de acciones inválidas \\
Verification depth & Verificación final / verificación por paso / verificación de doble capa & Tasa de acierto & Latencia \\
\bottomrule
\end{tabular}
\caption{Matriz base de experimentos de ablación para Agentes}
\end{table}

\subsection{Clasificación de fallos y método de atribución}
Sin una clasificación unificada de fallos, el experimento cae en una "magia negra del ajuste de parámetros". Se recomienda dividir los fallos en cuatro categorías: fallos de percepción, fallos de razonamiento, fallos de ejecución y fallos de verificación. Cada categoría de fallo corresponde a una acción de corrección distinta, lo que evita la ruta ineficiente de "agrandar el modelo ante cualquier problema".

\begin{importantbox}{La atribución de fallos debe ser "ejecutable"}
El resultado de la atribución debe poder mapearse directamente a una acción de ingeniería:
\begin{itemize}
    \item Fallo de percepción $\rightarrow$ reforzar el analizador sintáctico y la verificación del estado de la página;
    \item Fallo de razonamiento $\rightarrow$ ajustar la plantilla de reasoning y el disparo de la reflexión;
    \item Fallo de ejecución $\rightarrow$ optimizar el encapsulado de herramientas y la estrategia de reintento;
    \item Fallo de verificación $\rightarrow$ ampliar la cobertura del verificador y el umbral de error.
\end{itemize}
\end{importantbox}

\begin{warningbox}{Evitar el "sesgo de las muestras exitosas"}
Si solo se analizan las trayectorias exitosas, el sistema juzgará mal su propia robustez. Se debe priorizar el análisis de las muestras de fallo de alto costo (costo elevado, cadena larga, no recuperable), pues estas muestras determinan el límite superior de viabilidad en producción.
\end{warningbox}

\subsection{Ritmo de iteración en producción}
El curso no recomienda desplegar cambios grandes de una sola vez. La ruta más segura es "shadow mode $\rightarrow$ despliegue gradual con tráfico reducido $\rightarrow$ despliegue total". Cada etapa debe tener condiciones de reversión claras y un monitoreo continuo del SLO a nivel de tarea.

\begin{knowledgebox}{Recomendaciones de monitoreo en producción}
\begin{itemize}
    \item Umbral de caída de la tasa de éxito y estrategia de degradación automática;
    \item Umbral de aumento súbito de costo y compuerta de presupuesto;
    \item Umbral de tasa de error de herramientas y mecanismo de corte automático.
\end{itemize}
\end{knowledgebox}

\subsection{Resumen de la sección}
El objetivo del diseño experimental no es "demostrar que el sistema es bueno", sino "ubicar por qué el sistema es bueno o malo". Reproducible, atribuible y reversible son las tres condiciones mínimas del sistema experimental de un Agente.

